The computational comparisons hold primitive shocks common before applying the different repayment rules. A disaster clause and a GDP clause therefore face the same output realization within a paired experiment. Differences in proceeds arise from the contractual mapping and the face chosen against it, rather than from assigning one instrument a more favorable sample of disasters. Across seeds, independent draws reveal the numerical uncertainty of the estimated maximum and of its paired gain. The reported standard error of the mean describes averaging those independent experiments; it is smaller than the standard deviation of one simulation run, which is retained in the results data.

This distinction is particularly useful for comparing nearby entries across tables. Small differences in maximizing face can reflect sample fluctuations around a relatively flat peak, while proceeds change very little. A reported standard error of zero at the displayed precision is a rounded numerical quantity, except for identities such as the constant-payoff equivalence. It does not mean that the growth process or the true fiscal limit is known without uncertainty. The model's parameter choices and institutional assumptions are held fixed in these repetitions and require separate sensitivity analysis.

The deferral exercise uses a different numerical object. It solves a joint borrowing-and-survival system over an inherited-debt state using conditional growth quantiles. Its table therefore reports grid refinements and fixed-point residuals, rather than presenting its entries as twenty-seed estimates from the one-period optimizer. The comparison is conditional on the specified distribution and on the interpolation scheme. Reporting the method and its boundary checks separately keeps simulation noise, discretization sensitivity, and economic specification uncertainty distinct.
